\section{Free Modules Revisited}
\subsection{$R$-mod}
Recall our previous definition of an $R$-module. A module over $R$ is an abelian group $M$, together with an action of $R$ on $M$. The action of $r\in R$ on $m\in M$ is denoted $r m$, which is the notation for \emph{left}-modules. We are assuming $R$ is commutative, so this difference is immaterial. Since an action is a ring homomorphism $R \to \End_{\mathsf{Ab} }(M) $, this tells us that
\begin{itemize}
	\item $(r_1+r_2)m=r_1m+r_2m$;
	\item $1m=m$ and $(r_1r_2)m=r_1(r_2m)$;
	\item $r(m_1+m_2)=r m_1+r m_2$.
\end{itemize}

Modules over $R$ form the ring $\Rmod{R}$. We are going to attempt to uncover many properties of $\Rmod{R} $ and $R$ in this chapter. We will start by looking at the full subcategory of $\Rmod{R} $ composed of free $R$-modules and whos morphisms are $R$-linear group homomorphsisms.

\subsection{Linear Independence and Bases}
A free $R$-module on a set $S$, denoted $F^R(S)$, is an $R$-module containing $S$ which is initial with respect to the existence of a set map $j : S \to F^R(S)$. That is,
\[\begin{tikzcd}[row sep = 2.5em, column sep = 2.5em]
	F^R(S)\arrow[r, dotted, "\phi "]&M\\
	S\arrow[u, "j"]\arrow[ru, "i"]
\end{tikzcd}\]
In this case, $\phi $ is an $R$-linear map, and $i,j$ are set functions. We showed $F^R(S)$ can be explicitly realized as $R^{\oplus S} $.

Our main tool in the study of these modules will be linear independence and bases. Recall that for all sets $I$, there is a canonical injection $j : I \to F^R(I)$, and any function $i : I \to M$ determines a unique $R$-module homomorphism $\phi : F^R(I) \to M$ such that the diagram
\[\begin{tikzcd}[row sep = 2.5em, column sep = 2.5em]
	F^R(I)\arrow[r, "\phi "]&M\\
	I\arrow[u, "j"]\arrow[ru, "i"]
\end{tikzcd}\]
commutes. We can apply this universal property in the case that $I$ is an indexing set and $i$ is an indexing map in $M$.

\begin{definition}[Linear Independence, Generating Sets]\label{dfn:6.1}
	Let $i : I \to M$ be an indexing set. We say that $i$ is linearly independent if the canonical $R$-linear map $\phi : F^R(I) \to M$ is injective; $i$ is linearly dependent otherwise. We say that $i$ generates $M$ if $\phi $ is surjective.
\end{definition}

We generally refer not to the indexing map $i$, but instead the subset $i(I)$ of $M$ determined by the map. A messier, but equivalent and more common definition of linear independence and generating sets is as follows: a set $S=\left\{ m_\alpha  \right\}_{\alpha \in I}\subseteq M $ is linearly independent if and only if the existstence a collection $\left\{ r_\alpha  \right\}_{\alpha \in I} \subseteq R$ such that
\[
	\sum_{\alpha \in I} r_\alpha m_\alpha =0
\]
implies that $r_\alpha =0\forall \alpha $. The set $S$ is said to generate $M$ if and only if every element of $M$ can be written as
\[
	\sum_{\alpha \in I} r_\alpha m_\alpha 
\]
for some collection $\left\{ r_\alpha  \right\}_{\alpha \in I} \subseteq R$. Note here that only \emph{finite} linear combinations are defined, so if $I$ is infinite, we implicitly assume $m_\alpha =0$ for all but finitely many $\alpha $.

\begin{lemma}\label{lma:6.1}
	Let $M$ be an $R$-module, and let $S\subseteq M$ be a linearly independent subset. Then there exists a maximal linearly independent subset of $M$ containing $S$.
\end{lemma}
\begin{proof}
	Let $\mathscr{S} $ be the collection of linearly independent subsets of $M$ containing $S$, which can be ordered by inclusion. Since $S\subseteq S$, $\mathscr{S} \neq \varnothing $. By Zorn's lemma, it suffices to show that every chain in $\mathscr{S} $ has an upper bound. Apparently the union of a chain of linearly independent subsets is linearly independent, but I don't understand the reasoning. I'll come back to this.
\end{proof}

This motivates the following definition as well.
\begin{definition}[Basis]\label{dfn:6.2}
	A set $B\subseteq M$ (an indexed set $B \to M$) is called a basis of $M$ if and only if it is linearly independent and generates $M$.
\end{definition}

Clearly, any basis is a maximal linearly independent set, since any $m\notin B$ can be generated by $B$, and thus attempting to include it would make the map $F^R(B)\to M$ no longer injective. However, the converse is not true for general rings, and it will only be true for vector spaces. This is why modules over a field are so special. Recalling our intiial definition \ref{dfn:6.1}, we can prove the follwing lemma.
\begin{lemma}\label{lma:6.2}
	An $R$-module $M$ is free if and only if it admits a basis. Similarly, $B\subseteq M$ is a basis if and only if the natural homomorphism $R^{\oplus B} \to B$ is an isomorphism.
\end{lemma}
\begin{proof}
	This is immediate. By definition \ref{dfn:6.2}, $B$ must be linearly independent and generate $M$. By definition \ref{dfn:6.1}, we have that $\phi : F^R(B) \to M$ is injective and surjective. Conversely, if $\phi : F^R(B) \to M$ is an isomorphism, then $M$ admits $i(B)$ as a basis, for $i : B \to M$ the indexing map.
\end{proof}

Thus, the choice of a basis of $M$ amounts to the choice of an isomorphism $R^{\oplus B} \cong M$.

\subsection{Vector Spaces}
Using what we have already seen, we can prove the fundamental result that all modules over a field are free. Modules over a field $k$ are called $k$-vector spaces, elements of vector spaces are called vectors, and elements of $k$ are called scalars.

\begin{lemma}\label{lma:6.3}
	Let $k$ be a field, and let $V$ be a $k$-vector space. Let $B$ be a maximal linearly independent subset of $V$. Then $B$ is a basis of $V$.
\end{lemma}
\begin{proof}
	Let $v\in V\setminus B$. Clearly $B\cup \left\{ v \right\} $ is not linearly independent, by the maximality of $B$, and thus there exist $\left\{ c_{i} \right\}_{i=0}^t \subseteq k$, $\left\{ b_i \right\}_{i=1}^t\subseteq B$ such that
	\[
		c_0v+\sum_{i=1}^{t} c_ib_i=0
	.\]
	Note that if $c_0=0$, then we would have $c_i\neq 0$ for some $i\neq 0$, and thus $B$ would not be linearly independent. We thus have that $c_0$ is a unit, because $k$ is a field, and thus
	\[
		v=-\sum_{i=1}^{t} c_0^{-1} c_ib_i
	.\]
	Therefore, $B$ generates $V$.
\end{proof}
\begin{proposition}\label{prp:6.1}
	Let $k$ be a field, and $V$ a $k$-vector space. Let $S$ be a linearly independent set of vectors in $V$. Then there exists a basis of $V$ containing $S$.
\end{proposition}
\begin{proof}
	By lemma \ref{lma:6.1}, there exists a maximal linearly independent subset of $M$ containing $S$. By lemma \ref{lma:6.3}, since $k$ is a field, $B$ is a basis. Moreover, by lemma \ref{lma:6.2}, $V$ is free.
\end{proof}
\begin{lemma}\label{lma:6.4}
	Let $k$ be a field, and let $V$ be a $k$-vector space. Let $B$ be a minimal generating set for $V$. Then $B$ is a basis for $V$.
\end{lemma}
\begin{proof}
	Done in the exercises.
\end{proof}

Lemmas \ref{lma:6.3} and \ref{lma:6.4} fail over more general rings. However, putting them together, we find that a subset $B$ of a vector space is a basis $\iff $ $B$ is a minimal generating set $\iff$ $B$ is a maximal linearly independent set. This also shows all modules over a field are necessarily free, since any nonzero singleton is necessarily linearly independent, so a linearly independent subset of $V$ always exists.

\subsection{Recovering $B$ from $F^R(B)$}
We will now present a reconstruction of $B$ up to bijection. This result justifies the notion of dimension of a vector space and rank of a module.
\begin{proposition}\label{prp:6.2}
	Let $R$ be an integral domain, and let $M$ be a free $R$-module. Let $B$ be a maximally linearly independent subset of $M$, and let $S$ be a linearly independent subset of $M$. Then $\left| S \right| \leq \left| B \right| $. In particular, any two maximally linearly independent subsets of a free module over an integral domain have the same cardinality.
\end{proposition}
\begin{proof}
	Exercise 1.7 reduces the problem to proving the statement for $k$-vector spaces. It suffices to show there is an injective map $j : S \hookrightarrow B$, which by the canonical decomposition for functions of sets can be viewed as a bijection and an inclusion. Let $\leq $ be a well ordering on $S$, which exists by the Axiom of Choice, let $v\in S$, and suppose $j$ is defined for all $w\in S$ such that $w<v$. Suppose inductively that it is possible to define $j : S \to B$ for all $w<v$ such that the set $B'$, obtained by replacing $j(w)$ with $w$, is maximally linearly independent. We wish to show we can define $j(v)$ such that replacing $j(v)$ with $v$ will maintain the maximal linear independence of $B'$. Moreover, we want $j(w)\neq j(v)$ for all $w<v$ (injectivity for the lesser elements holds by inductive hypothesis). Transfinite induction will then show that $j$ is defined and injective on $S$.

	Now to show $j(v)$ may be defined as such. Consider $B'\cup \left\{ v \right\} $ as an indexed set. Since we consider it as an indexed set, even if $v\in B'$, we still have that $B'\cup \left\{ v \right\} $ is linearly dependent, as there exists a linear dependence relation
	\[
		c_0v + c_1b_1+\cdots +c_tb_t = 0
	\]
	with $b_i\in B$ distinct for all $i$ and $c_0\neq 0$. Since this is a linear dependence, there exists at least one $c_i\neq 0$ with $i\neq 0$, and thus we can assume WLOG that $c_1\neq 0$ and $b_1\in B'\setminus S$, since $S$ was linearly independent and $v\in S$. Therefore, $b_1\neq w$ for all $w<v$. Set $j(v)\coloneqq b_1$. Note that we have
	\[
		v=-c_0^{-1} \sum_{i=1}^t c_ib_i
	.\]
	It suffices to show replacing $b_1$ with $v$ maintains the linear independence of this set. Define $B''$ as $B'$ but with $b_1$ replaced by $v$. Note that any linear combination in $B''$ can be written as a linear combination in $B'$ by the last equation, and thus $B''$ is linearly independent. Since $B'$ is maximally linearly independent, so is $B''$.
\end{proof}
\begin{corollary}\label{cor:6.1}
	Let $R$ be an integral domain, and let $A,B$ be sets. Then
	\[
		F^R(A)\cong_{\Rmod{R} } F^R(B)\iff A\cong_{\mathsf{Set} }  B
	.\]
\end{corollary}
\begin{proof}
	Proven in the exercises.
\end{proof}

This shows that integral domains satisfy the ``Invariant Basis Number'' (IBN) property; that is, $R^m\cong R^n$ if and only if $n=m$. This holds over commutative rings as well, but not all rings satisfy the IBN property.

One way to interpret this is that the category of finitely generated free modules over an integral domain is classified by $\mathbb{Z} ^{\geq 0} $ up to isomorphism, as there is precisely one free module per nonnegative integer. This motivates the following definitions.
\begin{definition}[Rank/Dimension]\label{dfn:6.3}
	Let $R$ be an integral domain. The rank of a free $R$-module $M$, denoted $\rank _R M$, is the cardinality of a maximal linearly independent subset of $M$. The rank of a vector space $V$ is called the dimension, denoted $\dim_k V$.
\end{definition}

We generally omit the subscripts $R$, $k$ if context allows it. However, it's important to still be precise about what we are describing, since for example, $\dim_{\mathbb{R} } \mathbb{C} =2$ but $\dim_\mathbb{C} \mathbb{C} =1$.

\begin{proposition}\label{prp:6.3}
	Let $R$ be an integral domain, and let $M$ be a free $R$-module. Assume that $M=\left\langle S \right\rangle $; that is, $M$ is generated by $S$. Then $S$ contains a maximally linearly independent subset of $M$.
\end{proposition}
\begin{proof}
	Trivial.
\end{proof}
